% Generated draft: requires mathematical and source review.
\documentclass[11pt]{article}
\usepackage[margin=25mm]{geometry}
\usepackage{fontspec}
\setmainfont{Latin Modern Roman}
\usepackage{amsmath,amssymb,mathtools,unicode-math}
\setmathfont{Latin Modern Math}
\usepackage{enumitem,needspace,xurl,hyperref}
\setlength{\parindent}{0pt}
\setlength{\parskip}{5pt}
\setlength{\emergencystretch}{3em}
\newcommand{\N}{\mathbb{N}}
\newcommand{\Z}{\mathbb{Z}}
\newcommand{\Q}{\mathbb{Q}}
\newcommand{\R}{\mathbb{R}}
\newcommand{\C}{\mathbb{C}}
\newcommand{\F}{\mathbb{F}}
\newcommand{\A}{\mathbb{A}}
\newcommand{\PP}{\mathbb{P}}
\newcommand{\E}{\mathbb{E}}
\newcommand{\eps}{\varepsilon}
\newcommand{\dd}{\,\mathrm{d}}
\begin{document}
\section*{UnsolvedMath ID: 11000073}
\textbf{Source problem code:} catalogue entry
\subsection{Definitions and mathematical statement}
Put \(G=\operatorname{Mod}_{g,1}\) in the source's convention, and let
\[1\longrightarrow Z\longrightarrow\widetilde G\longrightarrow G\longrightarrow1\]
be its stated universal central extension. A central extension has kernel contained in the center; universality means every central extension of \(G\) receives a unique compatible homomorphism from \(\widetilde G\). In the usual perfect-group setting its kernel identifies with \(H_2(G;\Z)\); the source does not specify the genus range or whether the one distinguished feature is a puncture, marked point or boundary component. These conventions and perfectness must be fixed before naming this extension. No cyclic-kernel description is imposed.

A group \(A\) is residually finite if
\[\forall a\in A\setminus\{1\},\ \exists\text{ finite group }F,\ \exists\phi:A\to F,\quad\phi(a)\ne1.\]
The question asks whether \(\widetilde{\operatorname{Mod}}_{g,1}\) is residually finite, or instead has a nonidentity element invisible in every finite quotient. The source recalls Grossman's residual finiteness of \(\operatorname{Mod}_g\) and \(\operatorname{Mod}_{g,1}\), but that of the quotient alone does not answer the central-extension question. Given that quotient property, elements outside \(Z\) are already separated through \(G\), so finite detection of the remaining nontrivial central-kernel elements is the essential additional requirement.

The source also recalls residual finiteness of \(\operatorname{Sp}(2g,\Z)\), with
\[\bigcap_{N\geq1}\ker\bigl(\operatorname{Sp}(2g,\Z)\to\operatorname{Sp}(2g,\Z/N\Z)\bigr)=\{1\}.\]
This is context about that integral matrix group, not the requested extension.

The research summary's statement that finitely generated subgroups of a residually finite mapping class group need not be residually finite is incompatible with the definition: restriction of a separating finite quotient separates any nonidentity subgroup element. Residual finiteness passes to every subgroup. That summary claim is therefore flagged rather than used to reinterpret the question as a subgroup problem.
\subsection{Short English statement}
Determine residual finiteness of the specified universal central extension of \(\operatorname{Mod}_{g,1}\). Known residual finiteness of the base group does not settle the extension, and the subgroup assertion in the research summary is erroneous.
\subsection{Sources}
\begin{itemize}
\item Local UnsolvedMath registry, permanent ID \(11000073\), exact source code \texttt{catalogue entry}.
\item Farb, \emph{Problems on Mapping Class Groups and Related Topics}, Question 6.4, PDF page 55: \url{https://www.math.uchicago.edu/~farb/papers/mcgbook.pdf}. Registry access 2026-07-29; pdf-font-structure extraction.
\item Grossman is credited with residual finiteness of the base mapping class groups; the summary dates this 1974/75 but supplies no full bibliographic entry.
\end{itemize}
\end{document}
